\documentclass[11pt]{article}
\usepackage[margin=24mm]{geometry}
\usepackage{amsmath,amssymb,amsthm}
\newtheorem{theorem}{Theorem}
\newtheorem{lemma}[theorem]{Lemma}
\title{GST whole-space causal folds: all-dimensional Hodge spectra, fiber multiplicity, and topology-changing realizations}
\date{21 October 2026}
\begin{document}
\maketitle

\section{Exact hypothesis and realized space types}
GST GeneralSpace has no intrinsic complex dimension or Hodge structure.
A realized world has additional data: a geometric space, a complex
or mixed-Hodge realization, a graded cohomology theory, and (where
appropriate) an algebraic-cycle class map. In a GST self-expanding
cosmos each newborn world may have DIFFERENT internal state and law
types. A causal graph edge does not automatically induce a Hodge
morphism. We must separately specify its actual geometric transport.

The user's experiment assumes the rational Hodge conjecture TRUE for
all smooth projective worlds and asks whether folding the underlying
space, its dimension, and its causal realization yields a universal
Hodge structure or a different shape. We test the ENTIRE Hodge
decomposition, without initially selecting rational cohomology
classes or claiming a new Hodge conjecture proof.

\section{The precise multiplicity hidden by cosmic-diagonal folding}
The actual Lean file GSTClassicalHodgeFiberedCosmology uses a
weight-p rational Hodge basis index i and sends every address
(p,i) to one common GST diagonal address (p,p) under
forgetMultiplicityToGST.

\begin{theorem}[Multiplicity-loss and indistinguishable base folds]
For a true rational Hodge fiber of rank r greater than zero,
the restriction of multiplicity forgetting is
\[
 F_p:\mathbb Q^r\longrightarrow\mathbb Q,\qquad
 (a_1,\ldots,a_r)\mapsto\sum_i a_i.
\]
Consequently its kernel has dimension r-1. Two rational
fiber-fold maps may induce IDENTICAL maps on the collapsed
cosmic address but have DIFFERENT ranks.
\end{theorem}
\begin{proof}
The Lean theorem forgetMultiplicityToGST_single sends every
basis vector to the same unit generator; linearity and
rank-nullity give the claim. For an explicit witness, let
\[
 M_A=\begin{pmatrix}
1&0&0\\0&1&0\\0&0&1\\0&0&0
\end{pmatrix},\qquad
 M_B=\begin{pmatrix}
1&1&1\\0&0&0\\0&0&0\\0&0&0
\end{pmatrix}.
\]
Both satisfy (1,1,1,1)M=(1,1,1), yet their ranks are
3 and 1. Therefore target cokernels have dimensions
1 and 3. These are exact RATIONAL LINEAR-ALGEBRA
discriminators, not claimed geometric correspondences.
\end{proof}

\section{All-N spectral law for the WHOLE Hodge structure}
Let X be any smooth projective complex N-fold. Its genuine
Hodge decomposition is
\[
 H^m(X,\mathbb C)=\bigoplus_{a+b=m}H^{a,b}(X).
\]
Define the real Hodge circle infinitesimal generator G
by its complex eigenvalues
\[
 \boxed{G|_{H^{a,b}}=i(a-b)\operatorname{id}.}       \tag{1}
\]
This whole operator belongs to the actual geometric
REALIZATION and is not specified by the dimensionless
GST carrier or by its abstract charge lattice alone.

\begin{theorem}[Universal polynomial for the emerging full diagonal]
For every integer N at least zero and each 0<=p<=N,
the operator
\[
 \boxed{
 P_{X,p}=\prod_{r=1}^{p}
       \left(\operatorname{id}
                 +\frac{G^2}{4r^2}\right)
            \quad\text{on }H^{2p}(X,\mathbb C)          } \tag{2}
\]
is idempotent, and its image is precisely the ENTIRE
complex Hodge diagonal \(H^{p,p}(X)\). Equivalently,
\[
 \boxed{H^{p,p}(X)=\ker(G|_{H^{2p}(X,\mathbb C)}).}   \tag{3}
\]
For a projective complex torus of dimension N,
\[
 h^{a,b}={N\choose a}{N\choose b},\qquad
 \operatorname{rank}P_{X,p}={N\choose p}^{2}.       \tag{4}
\]
This complex dimension is distinct from the
dimension over rational numbers of the rational
Hodge classes in the diagonal.
\end{theorem}
\begin{proof}
On H^{a,b} in degree 2p, the signed charge
a-b is in the set \(-2p,-2p+2,\ldots,2p\).
If a-b=0, every polynomial factor is one.
If a-b=+/-2s for 1<=s<=p, the factor with r=s
is zero. Hence (2) is exactly the Hodge-diagonal
projector and squares to itself for ALL N,p.
The torus formula follows from the exterior
algebra on the N holomorphic and N anti-holomorphic
one-forms. No conjecture about algebraic cycles
is used in this theorem.
\end{proof}

\section{Causal shape curvature and Hodge birth/loss dimensions}
For a typed GST causal edge from realized world X to
realized world Y with an ACTUAL certified cohomological
transport \(T:H^{2p}(X,\mathbb C)\to H^{2q}(Y,\mathbb C)\),
possibly with the correct Tate shift, define
\[
 \boxed{\mathcal K_T=G_YT-TG_X.}                   \tag{5}
\]
For two composable realized edges T and U,
\[
 \boxed{\mathcal K_{UT}=\mathcal K_U T
                              +U\mathcal K_T.}  \tag{6}
\]
This is the EXACT causal cocycle law for changing
whole Hodge REALIZATIONS, not a generic coordinate
identity on a fixed original X.

\begin{theorem}[Whole-shape causality criterion]
The certified transport T intertwines EVERY Hodge
charge iff \(\mathcal K_T=0\). The old diagonal is
taken into the new diagonal iff
\[
 \boxed{(I-P_{Y,q})T P_{X,p}=0.}                   \tag{7}
\]
If this weaker condition holds, define
\(T^{\mathrm{diag}}:H^{p,p}(X)\to H^{q,q}(Y)\).
Its target birth and source loss dimensions are
\[
 \boxed{
 B_T=h^{q,q}(Y)-\operatorname{rank}T^{\mathrm{diag}},
 \quad L_T=h^{p,p}(X)-\operatorname{rank}T^{\mathrm{diag}},
 \quad B_T-L_T=h^{q,q}(Y)-h^{p,p}(X).}             \tag{8}
\]
A GST higher-causality 2-cell can certify equivalence
of two such events ONLY after an equality of their
realized cohomology maps has been provided.
\end{theorem}
\begin{proof}
Both Hodge generators are diagonalizable with
charge-labelled eigenvalues, so intertwiners are
exactly charge-preserving maps. Formula (7)
is the whole-diagonal projector criterion.
Formula (6) expands and cancels the intermediate
world generator, and (8) is rank-nullity.
The statement needs an actual T; merely having a
globular higher-causal identity cell does not supply it.
\end{proof}

\section{An actual varying-complex-shape ALL-N test; no nonholomorphic action on a selected class}
For every N>=2, take genuine projective abelian worlds
\[
 X_{N,t}=E_i\times E_{i+t}\times E_i^{N-2},
 \qquad t\in\mathbb R,\ |t|\ll1,
\]
where E_tau is the complex elliptic torus with period tau.
The smooth rational cohomology local system is marked by
\(e_j=dx_j,f_j=dy_j\). The ACTUAL complex structure
changes with t, and hence changes its whole Hodge generator G_t.

Put \(\eta=e_1f_2-f_1e_2=e_1f_2+e_2f_1\).
The first two holomorphic forms are
\(\omega_1=e_1+if_1,\omega_2=e_2+(i+t)f_2\).
The exact period equation is
\[
 \boxed{\eta\wedge\omega_1\wedge\omega_2
        =t\,e_1e_2f_1f_2.}                       \tag{9}
\]
At t=0 eta is the class of an ACTUAL algebraic
combination of the diagonal and the two coordinate
fibers. At t!=0 eta is NOT type (1,1).

Even more directly, in the ordered exterior basis
\((e_1e_2,e_1f_1,e_1f_2,e_2f_1,e_2f_2,f_1f_2)\)
the actual WHOLE-SHAPE Hodge generator gives
\[
 \boxed{
 G_t^{(2)}\eta=t(e_1f_2-e_2f_1)+t^2f_1f_2,
 \qquad G_0^{(2)}\eta=0.}                        \tag{10}
\]
This calculation is from the changing holomorphic
basis and its conjugate, not a fixed-complex-structure
nonholomorphic pullback.
For all 1<=p<N, the marked class
\[
 \alpha_{N,p}=\eta\prod_{j=3}^{p+1}(e_jf_j)
\]
is (p,p) at t=0, and has nonzero off-diagonal
(p+1,p-1) and (p-1,p+1) components for t!=0,
by the ACTUAL Kunneth decomposition.
The Hodge NUMBERS of a complex torus stay constant
while the geometric Hodge SUBSPACES vary: this is an
important distinction from birth/loss of Betti dimensions.
The hypothesis of universal Hodge algebraicity is
consistent with every realized fiber in this family.

\section{An ACTUAL topology-changing ALL-N causal fold with Hodge polynomial}
For every N>=1, take the projective quadratic family
\[
 Q_t^N=\{z_0^2+\cdots+z_N^2+t w^2=0\}
                      \subset\mathbb P^{N+1}.
\]
For t!=0 the world is smooth projective of dimension N.
At t=0 it becomes the actual singular cone over a smooth
projective quadric \(Q^{N-1}\). Thus this changes the
underlying topology, not merely the Hodge coordinates.

Let \(E(T;u,v)\) denote the Hodge--Deligne polynomial,
and put \(w=uv\).
The smooth quadric's algebraic cell decomposition gives
\[
 E(Q^N_{\rm sm})=\sum_{j=0}^N w^j+
  \begin{cases}w^{N/2}&N\text{ even}\\0&N\text{ odd.}\end{cases}
\]
The singular cone minus its vertex is an
\(\mathbb A^1\)-bundle over \(Q^{N-1}\), so
\(E(Q^N_{\rm cone})=1+wE(Q^{N-1}_{\rm sm})\).

\begin{theorem}[An EXACT all-N Hodge-shape jump across a true topology fold]
For every N>=1,
\[
 \boxed{
 E(Q^N_{\rm sm};u,v)-E(Q^N_{\rm cone};u,v)
      =(-1)^N(uv)^{\lceil N/2\rceil}.}           \tag{11}
\]
In particular,
\[
 \boxed{\chi(Q^N_{\rm sm})-\chi(Q^N_{\rm cone})
          =(-1)^N.}                              \tag{12}
\]
The local fold is an ordinary quadratic critical
point whose Milnor fiber has homotopy type S^N.
For even N=2p the smooth quadric has an EXTRA
middle (p,p) geometric direction: two actual
rulings rather than one ambient class.
For odd N the virtual Hodge--Deligne correction
has the opposite sign. The singular wall must
be interpreted with MIXED Hodge cohomology;
(11) is a polynomial identity, not an identification
of every individual mixed cohomology group.
\end{theorem}
\begin{proof}
A smooth complex quadric has one algebraic
affine cell at each dimension 0..N and an
additional middle cell exactly when N is even.
The singular cone stratifies into its one vertex
and an affine line bundle over the lower quadric.
Additivity of compact-support Hodge--Deligne
polynomials yields the two displayed formulas.
For N=2p their difference is +w^p. For
N=2p+1 it is -w^{p+1}. This proves (11);
setting u=v=1 proves (12). In the chart w=1,
the local smoothing equation \(\sum z_j^2+t=0\)
has one complex Morse critical point and an
N-sphere Milnor fiber, independently certifying
the parity of the Euler jump.
\end{proof}

\section{Interpretation UNDER the user's all-varieties-Hodge-TRUE assumption}
The full realized shape is the family of
spectral subspaces of G, together with the
causal maps and possible mixed-Hodge boundary
data. A newly born rational class in its
diagonal eigenspace is algebraic UNDER
the assumed Hodge conjecture.
A newly born class outside that eigenspace
is not a Hodge conjecture counterexample.
A possible genuine contradiction would
require BOTH a rational zero-charge class
AND a rigorous proof that it lies OUTSIDE
the image of the canonical Chow class map
on a smooth projective variety. The
shape eigen-equation alone decides the
first of those conditions, not the second.

The important outcome is NOT another
projective bundle or branch-cover packet:
equations (2),(5)--(8),(10),(11) form a
typed whole-realization, all-N Hodge/causality
test, and the multiplicity theorem proves
why the collapsed cosmic diagonal loses
exactly the information needed to execute it.
The source modules have been selectively
read at their exact current Git revision;
the repository contains 701 Lean files,
not all of which have had full body review.
\end{document}
